% mazzi: 15,16
\chapter{Tecnica, tattica e didattica del futsal}
\label{cap:tecnicatattica}

\textbf{Tecnica e tattica per il preparatore fisico}: conoscerne i principi serve a dialogare
correttamente con l'allenatore, che è il responsabile tecnico, e a \textbf{trasferire alla partita}
il lavoro svolto in campo.

Questi principi non servono al preparatore per sostituirsi all'allenatore o per strutturare da sé
situazioni tattiche, ma per indirizzare il proprio lavoro in una direzione che comprenda anche gli
aspetti cognitivi e la tattica --- individuale e di squadra --- della disciplina.

% ---------------------------------------------------------------------------
\section{L'impostazione dell'allenamento}

\rubrica{L'allenamento cognitivo}
Il calcio a 5 sottopone i giocatori a stimoli costanti e che si ripetono con grande frequenza, da
tutti i punti di vista: di conseguenza l'allenamento deve basarsi sulla \textbf{componente
cognitiva}.

\rubrica{L'integrazione del lavoro con e senza palla}
\begin{itemize}
  \item \textbf{con la palla}:
    \begin{itemize}
      \item \textit{small sided games};
      \item esercitazioni ad alta intensità con scopo prettamente fisico --- per esempio
        l'allenamento della \textbf{repeated sprint ability}, capacità fondamentale per il giocatore
        di calcio a 5;
    \end{itemize}
  \item \textbf{senza palla}: i lavori per le capacità metaboliche --- aerobiche e anaerobiche --- e
    quelli prettamente neuromuscolari.
\end{itemize}
Il preparatore alterna abitualmente le due famiglie di esercitazioni.

% ---------------------------------------------------------------------------
\section{L'organizzazione della fase di possesso}

Il discorso tattico si divide in \textbf{organizzazione della fase di possesso} --- quando la
squadra ha il pallone, cioè tutti i percorsi che lo portano verso la porta avversaria --- e
\textbf{organizzazione della fase di non possesso}.

\rubrica{Le due strutture di base}
\begin{itemize}
  \item il \textbf{3-1}: tre giocatori sfalsati --- uno nella zona più centrale del campo e due in
    ampiezza sulle fasce --- più un \textbf{pivot} in avanti, che nel gergo dà la profondità;
  \item il \textbf{4-0}: quattro giocatori disposti a garantire l'ampiezza del campo, che cercano la
    profondità attraverso movimenti di smarcamento e attacchi alla profondità, per portare la palla
    il più vicino possibile alla porta avversaria.
\end{itemize}
La sistemazione subisce comunque cambiamenti continui, imposti dalla dinamica di gioco: si arriva
per esempio a situazioni in cui il possessore ha un compagno a \textbf{sostegno difensivo} e due
compagni oltre la linea della palla, che garantiscono la profondità.

\rubrica{Le due letture del 5 contro 5}
Il possessore di palla, grazie al movimento corretto dei compagni, può sviluppare il gioco secondo
due letture:
\begin{enumerate}
  \item \textbf{ai lati e in avanti}: a destra o a sinistra per il gioco in ampiezza, in avanti per
    quello in profondità;
  \item \textbf{dietro e avanti}: un passaggio dietro di sé, verso un compagno che gli garantisce
    sostegno o sicurezza, più due appoggi davanti a sé.
\end{enumerate}

\rubrica{Che cosa si osserva nelle azioni di possesso}
\begin{itemize}
  \item nella \textbf{disposizione a tre} si ricerca la profondità attraverso il pivot, anche con il
    passaggio diretto;
  \item quando la difesa avversaria è schierata molto bassa, la fase di possesso si prolunga ed è
    più difficile trovare lo spazio per attaccare la porta;
  \item nella \textbf{disposizione a quattro} non c'è profondità iniziale: i quattro tendono ad
    allinearsi per poi attaccare la profondità che hanno liberato, con movimenti di smarcamento in
    profondità e in ampiezza.
\end{itemize}

\rubrica{L'uscita dalla pressione}
\begin{itemize}
  \item \textbf{che cos'è}: nel gergo \textit{uscita press}, il modo in cui si porta la palla nella
    metà campo avversaria quando gli avversari pressano;
  \item \textbf{come si costruisce}: sono schemi e movimenti predeterminati, codificati e studiati
    durante l'allenamento;
  \item \textbf{perché serve}: ad alto livello --- internazionale, ma anche nazionale di alta
    qualificazione --- la pressione viene molto spesso portata molto alta, e le squadre sono
    costrette a studiare un modo per risolverla;
  \item \textbf{che cosa la caratterizza}: l'elevata frequenza di passaggi, controlli, smarcamenti,
    attacchi della profondità e \textbf{contromovimenti}. Sono qualità di tattica individuale e
    tecniche che fanno la differenza fra il giocatore di alto livello e gli altri.
\end{itemize}

% ---------------------------------------------------------------------------
\section{Le situazioni di gioco}

\subsection{Le parità numeriche, da 1 contro 1 a 3 contro 3}

\begin{center}
\begin{forest} classalbero
[Parità numeriche
  [Duello 1 contro 1]
  [Combinazioni a 2]
  [Combinazioni a 3]
]
\end{forest}
\end{center}

\begin{itemize}
  \item \textbf{duello 1 contro 1}: su un campo di 40 $\times$ 20\,m si gioca 4 contro 4 di
    movimento, esclusi i portieri, e i duelli individuali sono quindi molto frequenti. Per questo
    sono molto curati in allenamento, sia sul piano tecnico dall'allenatore sia su quello fisico dal
    preparatore;
  \item \textbf{combinazioni a 2}: situazioni di gioco che coinvolgono due giocatori della stessa
    squadra;
  \item \textbf{combinazioni a 3}: le stesse, estese a tre giocatori.
\end{itemize}

\subsection{Le superiorità numeriche: le transizioni}

Nascono quasi sempre da un passaggio di possesso fra le due squadre --- palla recuperata con un
contrasto, un intercetto o un semplice passaggio sbagliato --- da cui il termine
\textbf{transizioni}. Le più frequenti:
\begin{itemize}
  \item con tre giocatori coinvolti in fase di possesso: \textbf{3 contro 0} (la più frequente), poi
    3 contro 1 e 3 contro 2;
  \item \textbf{1 contro il portiere}, in gergo il ``tu per tu'';
  \item \textbf{2 contro il portiere}, indicata come 2 contro 0, e 2 contro 1.
\end{itemize}
Sono situazioni particolarmente curate dagli allenatori e da tutto lo staff nell'allenamento
settimanale.

\subsection{Le combinazioni a 2 e 3 giocatori}

I principi con cui si risolvono:
\begin{enumerate}
  \item \textbf{dai e segui};
  \item \textbf{sovrapposizione in ampiezza};
  \item \textbf{sovrapposizione in profondità};
  \item \textbf{diagonale e parallela}: il taglio diagonale, e la giocata in parallela, che è
    generalmente lungo la fascia laterale e in perpendicolare;
  \item \textbf{rotazioni}: un susseguirsi di trasmissioni di palla e di movimenti volti a liberare
    l'uomo per un 1 contro 1;
  \item \textbf{cambi di posizione}.
\end{enumerate}

\rubrica{Che cosa si osserva nelle combinazioni}
\begin{itemize}
  \item \textbf{gioco a 2 con sovrapposizione in ampiezza};
  \item \textbf{dai e segui con incrocio su appoggio in posizione di pivot}: appoggio al pivot,
    incrocio, tiro in porta;
  \item \textbf{gioco a 2 dai e segui}, dove l'elemento decisivo è lo \textbf{smarcamento}, e in
    particolare il contromovimento --- la finta con cui il giocatore trova lo spazio e il vantaggio
    che gli permettono di calciare;
  \item \textbf{gioco a 3 con pivot centrale} (dai e segui, con sovrapposizione in ampiezza) e con
    pivot esterno, che va a prendere la palla sull'esterno e gioca in parallela;
  \item \textbf{parallela e diagonale in forma ripetuta}: per smarcarsi, cioè per allontanarsi il
    più possibile dall'avversario e ricevere palla $\rightarrow$ cambi di direzione e di senso
    continui, cambi di appoggio e di baricentro costanti. È il collegamento diretto con la match
    analysis dal punto di vista cinematico.
\end{itemize}

% ---------------------------------------------------------------------------
\section{I fondamentali tecnici specifici}

\rubrica{I sette fondamentali del calcio}
\begin{multicols}{2}
\begin{enumerate}
  \item \textbf{calciare la palla}: tiro in porta e passaggio, detto anche trasmissione;
  \item \textbf{ricezione}: lo stop, cioè il controllo della palla;
  \item \textbf{guida della palla};
  \item \textbf{colpo di testa};
  \item \textbf{contrasto};
  \item \textbf{rimessa laterale};
  \item \textbf{tecnica del portiere}.
\end{enumerate}
\end{multicols}
Nel calcio a 5 la rimessa laterale si esegue con i piedi e con la palla ferma a terra
(cap.~\ref{cap:introfutsal}).

\rubrica{I sei fondamentali specifici del futsal}
\begin{enumerate}
  \item \textbf{stop di pianta}: con la suola, cioè con la parte inferiore del piede; garantisce di
    fermare la palla e quindi di eseguire il gesto tecnico successivo nel modo più preciso
    possibile;
  \item \textbf{stop orientato di pianta}: è una scelta di tattica individuale --- il controllo di
    pianta non serve a fermare la palla ma a spostarla, per trovare nuove linee di passaggio o per
    smarcarsi dall'avversario;
  \item \textbf{guida della palla con la pianta};
  \item \textbf{trasmissione di pianta}: oltre che con tutte le parti del piede come nel calcio, la
    palla viene spesso trasmessa di pianta, e a volte all'indietro di suola per fintare e ingannare
    l'avversario;
  \item \textbf{trasmissione di esterno collo accompagnato}: serve ad alzare il pallone, quindi a
    fare passaggi lunghi e alti che tagliano le linee difensive;
  \item \textbf{calcio di punta-piede}: con la parte anteriore del piede. È la differenza più
    spettacolare e importante rispetto al calcio: porta a tiri molto potenti e veloci, difficili da
    intercettare e da gestire per il portiere.
\end{enumerate}

\rubrica{Che cosa si osserva nell'esecuzione}
\begin{itemize}
  \item \textbf{conduzione di pianta}: controllo di suola $\rightarrow$ il giocatore prende il tempo
    all'avversario diretto e trova lo spazio per calciare;
  \item giocare su una superficie stabile e piatta con un pallone a \textbf{rimbalzo controllato}
    facilita l'esecuzione di questi gesti;
  \item \textbf{controllo orientato}: serve a smarcarsi e a superare l'avversario per tirare in
    porta;
  \item trasmissione di esterno collo accompagnato: eseguita sia nel breve sia per servire una palla
    alta e lunga;
  \item \textbf{calcio di punta}: il \textbf{piede d'appoggio} non è nella posizione classica
    insegnata nella didattica del tiro nel calcio --- parallelo e accanto al pallone --- ma in una
    posizione completamente diversa.
\end{itemize}

% ---------------------------------------------------------------------------
\section{Dal gesto tecnico alla richiesta neuromuscolare}

\rubrica{Che cosa chiede una finta}
In un dribbling eseguito da un giocatore di altissimo livello, ciò che interessa al preparatore
fisico non è il gesto ma i \textbf{gradi articolari} in cui va espressa forza:
\begin{itemize}
  \item l'atteggiamento di finta porta a una condizione coordinativa di equilibrio estremamente
    instabile;
  \item appoggio forte sulla caviglia sinistra, con il ginocchio opposto molto basso;
  \item elevato \textbf{grado di torsione} a livello pelvico e del \textit{core};
  \item la finta coinvolge perfino lo \textbf{sguardo}.
\end{itemize}
Da qui la domanda operativa: quali mezzi di allenamento adottare durante la settimana per rendere il
giocatore performante in queste condizioni.

\rubrica{Le due qualità da allenare}
\begin{enumerate}
  \item la \textbf{forza esplosiva}, per la quale saranno proposti test di valutazione specifici
    nelle lezioni successive;
  \item la capacità di \textbf{reagire all'instabilità}, che poggia su una capacità coordinativa
    fondamentale: l'\textbf{equilibrio motorio}.
\end{enumerate}

\rubrica{I mezzi di allenamento per l'instabilità}
\begin{itemize}
  \item esercizi che prevedono l'instabilità, con \textbf{supporti propriocettivi}: supporti
    morbidi, tavole propriocettive;
  \item instabilità indotta da \textbf{impatti sul \textit{core}};
  \item ogni operatore ha però le proprie convinzioni, e il modo più efficace per riprodurre la
    situazione specifica del calcio a 5 resta lavorare su \textbf{superfici rigide}, cioè su quelle
    su cui il giocatore è abituato a giocare.
\end{itemize}

\rubrica{Il cambio appoggio}
È il gesto in cui l'equilibrio, capacità coordinativa, diventa gesto tecnico. Fa parte del bagaglio
dei giocatori di alta qualificazione.
\begin{enumerate}
  \item il giocatore controlla la palla di suola con un piede --- per esempio il sinistro --- mentre
    è in appoggio sull'altro;
  \item lascia andare la palla dal piede che ha controllato, facendo sì che quello diventi il
    \textbf{nuovo appoggio};
  \item contemporaneamente esegue una finta con il resto del corpo, per eludere la marcatura del
    diretto avversario.
\end{enumerate}
Richiede capacità di equilibrio, \textbf{stiffness} muscolare e forza esplosiva nella ripartenza
successiva.

\rubrica{L'accelerazione}
\begin{itemize}
  \item dopo il cambio appoggio serve una grande \textbf{capacità di accelerazione}, cioè di
    aumentare la propria velocità nel minor tempo possibile;
  \item nel calcio esistono moltissimi studi sull'accelerazione e sulla sua allenabilità; nel calcio
    a 5 la capacità viene indagata nella \textbf{match analysis};
  \item è fondamentale per la natura del gioco: fasi continue di attacco alla porta avversaria e di
    difesa della propria, che richiedono grande intensità sul piano fisico-motorio oltre che su
    quello cognitivo;
  \item accelerazione e controllo motorio facilitano a loro volta l'esecuzione del gesto tecnico e
    della tattica, individuale e collettiva.
\end{itemize}

% ---------------------------------------------------------------------------
\section{Le situazioni da calcio fermo}

Per falli laterali e calci d'angolo ogni allenatore ha il proprio bagaglio e le squadre creano
moltissimi schemi, studiati e applicati in campo durante la settimana: allenarli e strutturarli è un
fattore fondamentale della programmazione settimanale, perché portano spesso al gol o comunque ad
azioni pericolose.

\rubrica{Il calcio di rigore}
\begin{itemize}
  \item \textbf{rigore parato}: contro un tiro molto potente di collo interno, il portiere si sposta
    rimanendo in posizione eretta, per occupare il maggior volume possibile;
  \item \textbf{rigore realizzato}: eseguito di collo, con estrema potenza e con precisione.
\end{itemize}

\rubrica{Il tiro libero}
\begin{itemize}
  \item \textbf{esecuzione}: senza opposizione degli avversari, generalmente con un tiro diretto in
    porta;
  \item tecnica del portiere: posizione molto avanzata rispetto alla linea di porta; resta eretto
    fino all'ultimo per occupare il volume più ampio possibile e indurre in errore chi calcia, e
    solo allora effettua la parata;
  \item \textbf{vincolo regolamentare}: il portiere non può avanzare prima del tiro;
  \item la posizione avanzata non basta sempre: contro un tiro effettuato di interno piede il
    portiere può non riuscire a opporsi.
\end{itemize}

\rubrica{Il calcio d'angolo}
\begin{itemize}
  \item primo schema: un giocatore effettua un piccolo taglio davanti per far passare il passaggio
    $\rightarrow$ tiro da fuori, facilitato anche dalla scelta tattica della difesa avversaria;
  \item secondo schema: schema predeterminato che porta al gol con un passaggio diretto in area.
\end{itemize}

\rubrica{Il fallo laterale}
Battuta da posizione molto profonda, vicina a quella del calcio d'angolo: lo schema predeterminato
porta allo smarcamento di un giocatore al limite dell'area, che va a segnare in maniera diretta.

% ---------------------------------------------------------------------------
\section{Che cosa se ne ricava per il lavoro del preparatore}

Conoscere questi principi serve a costruire insieme all'allenatore esercitazioni di campo che:
\begin{enumerate}
  \item abbiano come obiettivo principale quello fisico-atletico --- la sollecitazione della
    componente metabolica o di quella neuromuscolare;
  \item sollecitino anche l'\textbf{aspetto cognitivo}, dentro le conoscenze tecnico-tattiche del
    giocatore e la filosofia di gioco dell'allenatore;
  \item contengano \textbf{regole} che inducano nel giocatore i comportamenti tecnici e tattici
    voluti dall'allenatore.
\end{enumerate}
È il modo più efficace per curare la performance del giocatore e prepararlo alla partita.
